\startchapter[title={Resultados}, reference=Capitulo_6]

Este capítulo presenta los resultados de las pruebas definidas en el
Capítulo~\in[Capitulo_5]. Primero se comparan los dos modelos lineales
(Sección~\in[sec:comparacion_lineales]) y después el modelo Split con las
metaheurísticas (Sección~\in[sec:comparacion_split_metaheuristicas]). Todos
los tiempos son de reloj, medidos en una misma notebook en las condiciones
descriptas en la Sección~\in[sec:ejecucion].

%----------------------------------------------------------------------------------------
\startsection[title={Comparación entre modelos lineales}, reference=sec:comparacion_lineales]

Esta sección compara el modelo monolítico con el Split en tres situaciones:
las pruebas funcionales, una misma instancia del perfil realista en almacenes
de tamaño creciente y las pruebas de estrés. La comparación sigue las reglas
de la Sección~\in[sec:comparacion].

\startsubsection[title={Mismo óptimo}]

De los 60 pares comparables de las pruebas funcionales, los 49 que los dos
modelos resuelven tienen el mismo valor objetivo, con una diferencia máxima
de $3\mathord{,}83 \cdot 10^{-6}$ (casos 13 y 14), dentro de la tolerancia de
SCIP, y los 11 que son infactibles lo son en los dos modelos. Esto confirma
en la práctica el argumento de exactitud de la Sección~\in[sec:modelo_split],
también en las pruebas de cadena de frío, de velocidades y de elección de la
ruta más corta.

Los dos modelos, sin embargo, no siempre llegan a la misma solución. En 28
pares devuelven la misma; en los otros 21, el monolítico devuelve un óptimo
alternativo, con otras rutas (13 pares), otra asignación (16) o ambas cosas.
Es lo esperable: como el objetivo no penaliza la longitud de la ruta, el
monolítico puede elegir cualquier ruta que no empeore las penalizaciones,
mientras que la Etapa~1 del Split siempre devuelve la ruta mínima. Las
asignaciones distintas, por su parte, son empates del objetivo que cada
resolución rompe de forma distinta. El Anexo~\in[Anexo_Resultados] detalla
cada par.

\stopsubsection

\startsubsection[title={Tiempos en las pruebas funcionales}]

El Cuadro~\in[tab:res_tiempos_lineales] resume los tiempos de los 49 pares
resueltos.

\startplacetable[reference=tab:res_tiempos_lineales, location=here, title={Tiempos de los 49 pares resueltos en las pruebas funcionales. El tiempo total incluye escribir los datos y compilar los modelos; el del resolvedor cubre solo la última resolución de SCIP, que en el Split es la de la Etapa~2.}]
	\bTABLE[split=yes, align=middle, frame=on, width=fit]
	\setupTABLE[row][1][background=color, backgroundcolor=darkgray, style=bold]
	\bTR \bTH Medida \eTH \bTH Monolítico (s) \eTH \bTH Split (s) \eTH \bTH Monolítico / Split \eTH \eTR
	\bTR \bTD Tiempo total: suma \eTD \bTD 140,3 \eTD \bTD 24,4 \eTD \bTD 5,75 \eTD \eTR
	\bTR \bTD Tiempo total: mediana \eTD \bTD 0,225 \eTD \bTD 0,309 \eTD \bTD 0,73 \eTD \eTR
	\bTR \bTD Tiempo total: máximo \eTD \bTD 43,5 \eTD \bTD 3,9 \eTD \bTD 11,27 \eTD \eTR
	\bTR \bTD Tiempo del resolvedor: suma \eTD \bTD 124,8 \eTD \bTD 7,2 \eTD \bTD 17,35 \eTD \eTR
	\bTR \bTD Tiempo del resolvedor: mediana \eTD \bTD 0,007 \eTD \bTD 0,003 \eTD \bTD 2,33 \eTD \eTR
	\eTABLE
\stopplacetable

En la suma, el monolítico tardó 5,75 veces más que el Split, pero en la
mediana la relación se invierte: el monolítico fue más rápido en 38 de los
49 pares, con entre 0,37 y 0,98 veces el tiempo del Split. Son instancias
chicas, en las que el resolvedor tarda milisegundos y el tiempo total lo
pone la compilación, y compilar un único modelo cuesta menos que compilar los
de las dos etapas. La suma la dominan las instancias de grilla, en las que
el monolítico llega a tardar 28,9 veces más (caso 29), como muestra la
Figura~\in[fig:res_lineales_dispersion]. En el tiempo del resolvedor la
diferencia es mayor, porque el monolítico resuelve un \ac{mip} con más
variables y restricciones que cualquiera de las dos etapas.

\placefigure[force][fig:res_lineales_dispersion]{Tiempo total de cada par resuelto en las pruebas funcionales: monolítico contra Split (escala logarítmica en los dos ejes). Por encima de la diagonal, el monolítico fue más lento.}
{\startMPcode
numeric w, h; w := 8cm; h := 8cm;
drawarrow (0,0)--(w+0.4cm,0) withpen pencircle scaled 1pt;
drawarrow (0,0)--(0,h+0.4cm) withpen pencircle scaled 1pt;
draw (0.000cm,0.000cm)--(0.000cm,-0.120cm) withpen pencircle scaled .5pt;
label.bot(textext("\tfx 0,1"), (0.000cm,-0.120cm));
draw (2.667cm,0.000cm)--(2.667cm,-0.120cm) withpen pencircle scaled .5pt;
label.bot(textext("\tfx 1"), (2.667cm,-0.120cm));
draw (5.333cm,0.000cm)--(5.333cm,-0.120cm) withpen pencircle scaled .5pt;
label.bot(textext("\tfx 10"), (5.333cm,-0.120cm));
draw (8.000cm,0.000cm)--(8.000cm,-0.120cm) withpen pencircle scaled .5pt;
label.bot(textext("\tfx 100"), (8.000cm,-0.120cm));
draw (0.000cm,0.000cm)--(-0.120cm,0.000cm) withpen pencircle scaled .5pt;
label.lft(textext("\tfx 0,1"), (-0.120cm,0.000cm));
draw (0.000cm,2.667cm)--(-0.120cm,2.667cm) withpen pencircle scaled .5pt;
label.lft(textext("\tfx 1"), (-0.120cm,2.667cm));
draw (0.000cm,5.333cm)--(-0.120cm,5.333cm) withpen pencircle scaled .5pt;
label.lft(textext("\tfx 10"), (-0.120cm,5.333cm));
draw (0.000cm,8.000cm)--(-0.120cm,8.000cm) withpen pencircle scaled .5pt;
label.lft(textext("\tfx 100"), (-0.120cm,8.000cm));
label.bot(textext("Split, tiempo total (s, escala log.)"), (w/2,-0.85cm));
label.lft(textext("Monolítico, tiempo total (s, escala log.)") rotated 90, (-1.25cm,h/2));
draw (0.000cm,0.000cm)--(8.000cm,8.000cm) withpen pencircle scaled 1.5pt withcolor \MPcolor{blueline} dashed evenly;
fill fullcircle scaled 4pt shifted (1.107cm,0.921cm) withcolor \MPcolor{darkblue};
fill fullcircle scaled 4pt shifted (1.149cm,0.946cm) withcolor \MPcolor{darkblue};
fill fullcircle scaled 4pt shifted (1.249cm,0.873cm) withcolor \MPcolor{darkblue};
fill fullcircle scaled 4pt shifted (1.273cm,0.866cm) withcolor \MPcolor{darkblue};
fill fullcircle scaled 4pt shifted (1.087cm,0.749cm) withcolor \MPcolor{darkblue};
fill fullcircle scaled 4pt shifted (1.346cm,1.368cm) withcolor \MPcolor{darkblue};
fill fullcircle scaled 4pt shifted (1.402cm,1.207cm) withcolor \MPcolor{darkblue};
fill fullcircle scaled 4pt shifted (1.254cm,0.823cm) withcolor \MPcolor{darkblue};
fill fullcircle scaled 4pt shifted (1.021cm,0.824cm) withcolor \MPcolor{darkblue};
fill fullcircle scaled 4pt shifted (1.229cm,0.803cm) withcolor \MPcolor{darkblue};
fill fullcircle scaled 4pt shifted (1.175cm,0.899cm) withcolor \MPcolor{darkblue};
fill fullcircle scaled 4pt shifted (4.232cm,7.037cm) withcolor \MPcolor{darkblue};
fill fullcircle scaled 4pt shifted (4.187cm,7.036cm) withcolor \MPcolor{darkblue};
fill fullcircle scaled 4pt shifted (1.202cm,1.020cm) withcolor \MPcolor{darkblue};
fill fullcircle scaled 4pt shifted (1.638cm,1.872cm) withcolor \MPcolor{darkblue};
fill fullcircle scaled 4pt shifted (1.329cm,1.058cm) withcolor \MPcolor{darkblue};
fill fullcircle scaled 4pt shifted (1.226cm,0.769cm) withcolor \MPcolor{darkblue};
fill fullcircle scaled 4pt shifted (1.119cm,1.096cm) withcolor \MPcolor{darkblue};
fill fullcircle scaled 4pt shifted (1.407cm,1.154cm) withcolor \MPcolor{darkblue};
fill fullcircle scaled 4pt shifted (1.050cm,0.855cm) withcolor \MPcolor{darkblue};
fill fullcircle scaled 4pt shifted (1.854cm,4.541cm) withcolor \MPcolor{darkblue};
fill fullcircle scaled 4pt shifted (1.099cm,0.987cm) withcolor \MPcolor{darkblue};
fill fullcircle scaled 4pt shifted (1.069cm,0.896cm) withcolor \MPcolor{darkblue};
fill fullcircle scaled 4pt shifted (1.463cm,1.097cm) withcolor \MPcolor{darkblue};
fill fullcircle scaled 4pt shifted (1.213cm,0.915cm) withcolor \MPcolor{darkblue};
fill fullcircle scaled 4pt shifted (1.872cm,0.779cm) withcolor \MPcolor{darkblue};
fill fullcircle scaled 4pt shifted (1.661cm,5.305cm) withcolor \MPcolor{darkblue};
fill fullcircle scaled 4pt shifted (1.683cm,4.536cm) withcolor \MPcolor{darkblue};
fill fullcircle scaled 4pt shifted (2.114cm,3.145cm) withcolor \MPcolor{darkblue};
fill fullcircle scaled 4pt shifted (2.199cm,4.671cm) withcolor \MPcolor{darkblue};
fill fullcircle scaled 4pt shifted (1.736cm,4.762cm) withcolor \MPcolor{darkblue};
fill fullcircle scaled 4pt shifted (2.151cm,1.866cm) withcolor \MPcolor{darkblue};
fill fullcircle scaled 4pt shifted (1.654cm,1.502cm) withcolor \MPcolor{darkblue};
fill fullcircle scaled 4pt shifted (2.061cm,2.022cm) withcolor \MPcolor{darkblue};
fill fullcircle scaled 4pt shifted (1.135cm,0.921cm) withcolor \MPcolor{darkblue};
fill fullcircle scaled 4pt shifted (1.493cm,1.074cm) withcolor \MPcolor{darkblue};
fill fullcircle scaled 4pt shifted (1.457cm,5.353cm) withcolor \MPcolor{darkblue};
fill fullcircle scaled 4pt shifted (1.226cm,0.938cm) withcolor \MPcolor{darkblue};
fill fullcircle scaled 4pt shifted (1.184cm,0.847cm) withcolor \MPcolor{darkblue};
fill fullcircle scaled 4pt shifted (1.307cm,0.843cm) withcolor \MPcolor{darkblue};
fill fullcircle scaled 4pt shifted (1.120cm,0.912cm) withcolor \MPcolor{darkblue};
fill fullcircle scaled 4pt shifted (1.947cm,0.943cm) withcolor \MPcolor{darkblue};
fill fullcircle scaled 4pt shifted (1.555cm,0.870cm) withcolor \MPcolor{darkblue};
fill fullcircle scaled 4pt shifted (2.033cm,0.882cm) withcolor \MPcolor{darkblue};
fill fullcircle scaled 4pt shifted (2.000cm,0.837cm) withcolor \MPcolor{darkblue};
fill fullcircle scaled 4pt shifted (1.099cm,0.860cm) withcolor \MPcolor{darkblue};
fill fullcircle scaled 4pt shifted (1.110cm,0.832cm) withcolor \MPcolor{darkblue};
fill fullcircle scaled 4pt shifted (0.944cm,0.921cm) withcolor \MPcolor{darkblue};
fill fullcircle scaled 4pt shifted (1.009cm,0.843cm) withcolor \MPcolor{darkblue};
fill fullcircle scaled 4pt shifted (0.600cm,7.850cm) withcolor \MPcolor{darkblue};
label.rt(textext("\tfx par resuelto"), (1.000cm,7.850cm));
draw (0.300cm,7.430cm)--(0.900cm,7.430cm) withpen pencircle scaled 1.5pt withcolor \MPcolor{blueline} dashed evenly;
label.rt(textext("\tfx igual tiempo"), (1.000cm,7.430cm));
\stopMPcode}

\stopsubsection

\startsubsection[title={Escala}]

Para ver qué pasa cuando la instancia crece, se resolvió con los dos modelos
una misma instancia del perfil realista (Cuadro~\in[tab:perfil_realista]) en
grillas de 3×3 a 7×7, con 1 a 6 lotes y un tope de 10 minutos. En 4×4 la
instancia es infactible y no se grafica. La Figura~\in[fig:res_lineales_escala]
muestra el resultado, y el Anexo~\in[Anexo_Relevamiento] presenta el mismo
estudio desde el punto de vista del relevamiento.

\placefigure[force][fig:res_lineales_escala]{Tiempo hasta probar el óptimo de una misma instancia del perfil realista, según la grilla y la cantidad de lotes, con un tope de 10 minutos (escala logarítmica, la misma en los dos paneles).}
{\startcombination[2*1]
{\startMPcode
numeric w, h; w := 6.4cm; h := 6.5cm;
drawarrow (0,0)--(w+0.4cm,0) withpen pencircle scaled 1pt;
drawarrow (0,0)--(0,h+0.4cm) withpen pencircle scaled 1pt;
draw (0.533cm,0.000cm)--(0.533cm,-0.120cm) withpen pencircle scaled .5pt;
label.bot(textext("\tfx 3×3"), (0.533cm,-0.120cm));
draw (1.867cm,0.000cm)--(1.867cm,-0.120cm) withpen pencircle scaled .5pt;
label.bot(textext("\tfx 4×4"), (1.867cm,-0.120cm));
draw (3.200cm,0.000cm)--(3.200cm,-0.120cm) withpen pencircle scaled .5pt;
label.bot(textext("\tfx 5×5"), (3.200cm,-0.120cm));
draw (4.533cm,0.000cm)--(4.533cm,-0.120cm) withpen pencircle scaled .5pt;
label.bot(textext("\tfx 6×6"), (4.533cm,-0.120cm));
draw (5.867cm,0.000cm)--(5.867cm,-0.120cm) withpen pencircle scaled .5pt;
label.bot(textext("\tfx 7×7"), (5.867cm,-0.120cm));
draw (0.000cm,0.000cm)--(-0.120cm,0.000cm) withpen pencircle scaled .5pt;
label.lft(textext("\tfx 0,1"), (-0.120cm,0.000cm));
draw (0.000cm,1.226cm)--(-0.120cm,1.226cm) withpen pencircle scaled .5pt;
label.lft(textext("\tfx 1"), (-0.120cm,1.226cm));
draw (0.000cm,2.452cm)--(-0.120cm,2.452cm) withpen pencircle scaled .5pt;
label.lft(textext("\tfx 10"), (-0.120cm,2.452cm));
draw (0.000cm,3.679cm)--(-0.120cm,3.679cm) withpen pencircle scaled .5pt;
label.lft(textext("\tfx 100"), (-0.120cm,3.679cm));
draw (0.000cm,4.905cm)--(-0.120cm,4.905cm) withpen pencircle scaled .5pt;
label.lft(textext("\tfx 1.000"), (-0.120cm,4.905cm));
label.bot(textext("Grilla"), (w/2,-0.85cm));
label.lft(textext("Tiempo hasta el óptimo (s, escala log.)") rotated 90, (-1.25cm,h/2));
draw (0.000cm,4.633cm)--(6.400cm,4.633cm) withpen pencircle scaled .5pt dashed evenly;
label.lft(textext("\tfxx tope: 10 min"), (6.300cm,4.433cm));
draw (0.533cm,0.369cm)--(3.200cm,0.585cm)--(4.533cm,1.277cm)--(5.867cm,1.323cm) withpen pencircle scaled 1.5pt withcolor \MPcolor{darkblue};
fill fullcircle scaled 4pt shifted (0.533cm,0.369cm) withcolor \MPcolor{darkblue};
fill fullcircle scaled 4pt shifted (3.200cm,0.585cm) withcolor \MPcolor{darkblue};
fill fullcircle scaled 4pt shifted (4.533cm,1.277cm) withcolor \MPcolor{darkblue};
fill fullcircle scaled 4pt shifted (5.867cm,1.323cm) withcolor \MPcolor{darkblue};
draw (0.533cm,1.846cm)--(3.200cm,1.937cm)--(4.533cm,2.094cm)--(5.867cm,2.436cm) withpen pencircle scaled 1.5pt withcolor \MPcolor{redline};
fill fullcircle scaled 4pt shifted (0.533cm,1.846cm) withcolor \MPcolor{redline};
fill fullcircle scaled 4pt shifted (3.200cm,1.937cm) withcolor \MPcolor{redline};
fill fullcircle scaled 4pt shifted (4.533cm,2.094cm) withcolor \MPcolor{redline};
fill fullcircle scaled 4pt shifted (5.867cm,2.436cm) withcolor \MPcolor{redline};
draw (0.533cm,2.072cm)--(3.200cm,2.247cm)--(4.533cm,2.360cm) withpen pencircle scaled 1.5pt withcolor \MPcolor{blueline};
fill fullcircle scaled 4pt shifted (0.533cm,2.072cm) withcolor \MPcolor{blueline};
fill fullcircle scaled 4pt shifted (3.200cm,2.247cm) withcolor \MPcolor{blueline};
fill fullcircle scaled 4pt shifted (4.533cm,2.360cm) withcolor \MPcolor{blueline};
draw (0.533cm,3.982cm)--(3.200cm,4.093cm) withpen pencircle scaled 1.5pt withcolor \MPcolor{greenline};
fill fullcircle scaled 4pt shifted (0.533cm,3.982cm) withcolor \MPcolor{greenline};
fill fullcircle scaled 4pt shifted (3.200cm,4.093cm) withcolor \MPcolor{greenline};
draw (0.300cm,6.400cm)--(0.900cm,6.400cm) withpen pencircle scaled 1.5pt withcolor \MPcolor{darkblue};
label.rt(textext("\tfx 2 lotes"), (1.000cm,6.400cm));
draw (0.300cm,5.980cm)--(0.900cm,5.980cm) withpen pencircle scaled 1.5pt withcolor \MPcolor{redline};
label.rt(textext("\tfx 3 lotes"), (1.000cm,5.980cm));
draw (0.300cm,5.560cm)--(0.900cm,5.560cm) withpen pencircle scaled 1.5pt withcolor \MPcolor{blueline};
label.rt(textext("\tfx 4 lotes"), (1.000cm,5.560cm));
draw (0.300cm,5.140cm)--(0.900cm,5.140cm) withpen pencircle scaled 1.5pt withcolor \MPcolor{greenline};
label.rt(textext("\tfx 6 lotes"), (1.000cm,5.140cm));
\stopMPcode}{Split}
{\startMPcode
numeric w, h; w := 6.4cm; h := 6.5cm;
drawarrow (0,0)--(w+0.4cm,0) withpen pencircle scaled 1pt;
drawarrow (0,0)--(0,h+0.4cm) withpen pencircle scaled 1pt;
draw (0.533cm,0.000cm)--(0.533cm,-0.120cm) withpen pencircle scaled .5pt;
label.bot(textext("\tfx 3×3"), (0.533cm,-0.120cm));
draw (1.867cm,0.000cm)--(1.867cm,-0.120cm) withpen pencircle scaled .5pt;
label.bot(textext("\tfx 4×4"), (1.867cm,-0.120cm));
draw (3.200cm,0.000cm)--(3.200cm,-0.120cm) withpen pencircle scaled .5pt;
label.bot(textext("\tfx 5×5"), (3.200cm,-0.120cm));
draw (4.533cm,0.000cm)--(4.533cm,-0.120cm) withpen pencircle scaled .5pt;
label.bot(textext("\tfx 6×6"), (4.533cm,-0.120cm));
draw (5.867cm,0.000cm)--(5.867cm,-0.120cm) withpen pencircle scaled .5pt;
label.bot(textext("\tfx 7×7"), (5.867cm,-0.120cm));
draw (0.000cm,0.000cm)--(-0.120cm,0.000cm) withpen pencircle scaled .5pt;
label.lft(textext("\tfx 0,1"), (-0.120cm,0.000cm));
draw (0.000cm,1.226cm)--(-0.120cm,1.226cm) withpen pencircle scaled .5pt;
label.lft(textext("\tfx 1"), (-0.120cm,1.226cm));
draw (0.000cm,2.452cm)--(-0.120cm,2.452cm) withpen pencircle scaled .5pt;
label.lft(textext("\tfx 10"), (-0.120cm,2.452cm));
draw (0.000cm,3.679cm)--(-0.120cm,3.679cm) withpen pencircle scaled .5pt;
label.lft(textext("\tfx 100"), (-0.120cm,3.679cm));
draw (0.000cm,4.905cm)--(-0.120cm,4.905cm) withpen pencircle scaled .5pt;
label.lft(textext("\tfx 1.000"), (-0.120cm,4.905cm));
label.bot(textext("Grilla"), (w/2,-0.85cm));
label.lft(textext("Tiempo hasta el óptimo (s, escala log.)") rotated 90, (-1.25cm,h/2));
draw (0.000cm,4.633cm)--(6.400cm,4.633cm) withpen pencircle scaled .5pt dashed evenly;
label.rt(textext("\tfxx tope: 10 min"), (0.100cm,4.833cm));
draw (0.533cm,0.585cm)--(3.200cm,1.862cm)--(4.533cm,2.612cm)--(5.867cm,2.762cm) withpen pencircle scaled 1.5pt withcolor \MPcolor{darkblue};
fill fullcircle scaled 4pt shifted (0.533cm,0.585cm) withcolor \MPcolor{darkblue};
fill fullcircle scaled 4pt shifted (3.200cm,1.862cm) withcolor \MPcolor{darkblue};
fill fullcircle scaled 4pt shifted (4.533cm,2.612cm) withcolor \MPcolor{darkblue};
fill fullcircle scaled 4pt shifted (5.867cm,2.762cm) withcolor \MPcolor{darkblue};
draw (0.533cm,2.050cm)--(3.200cm,2.372cm)--(4.533cm,3.853cm)--(5.867cm,4.633cm) withpen pencircle scaled 1.5pt withcolor \MPcolor{redline};
fill fullcircle scaled 4pt shifted (0.533cm,2.050cm) withcolor \MPcolor{redline};
fill fullcircle scaled 4pt shifted (3.200cm,2.372cm) withcolor \MPcolor{redline};
fill fullcircle scaled 4pt shifted (4.533cm,3.853cm) withcolor \MPcolor{redline};
draw fullcircle scaled 6pt shifted (5.867cm,4.633cm) withpen pencircle scaled 1pt withcolor \MPcolor{redline};
draw (0.533cm,2.231cm)--(3.200cm,3.278cm)--(4.533cm,4.633cm) withpen pencircle scaled 1.5pt withcolor \MPcolor{blueline};
fill fullcircle scaled 4pt shifted (0.533cm,2.231cm) withcolor \MPcolor{blueline};
fill fullcircle scaled 4pt shifted (3.200cm,3.278cm) withcolor \MPcolor{blueline};
draw fullcircle scaled 6pt shifted (4.533cm,4.633cm) withpen pencircle scaled 1pt withcolor \MPcolor{blueline};
draw (0.533cm,4.026cm)--(3.200cm,4.633cm) withpen pencircle scaled 1.5pt withcolor \MPcolor{greenline};
fill fullcircle scaled 4pt shifted (0.533cm,4.026cm) withcolor \MPcolor{greenline};
draw fullcircle scaled 6pt shifted (3.200cm,4.633cm) withpen pencircle scaled 1pt withcolor \MPcolor{greenline};
draw fullcircle scaled 6pt shifted (0.600cm,6.400cm) withpen pencircle scaled 1pt withcolor \MPcolor{darkblue};
label.rt(textext("\tfx no termina (cortado al tope)"), (1.000cm,6.400cm));
\stopMPcode}{Monolítico}
\stopcombination}

El contraste es de forma y no de grado. En el Split las curvas crecen
despacio: con hasta 4 lotes resuelve cualquier grilla en menos de 10
segundos. En el monolítico, en cambio, se disparan: con 3 lotes pasa de 9
segundos en 5×5 a 139 en 6×6, y no prueba el óptimo dentro del tope con 3
lotes en 7×7, 4 lotes en 6×6 y 6 lotes en 5×5. Entre los puntos que los dos
cierran, el monolítico tarda hasta 27 veces más (3 lotes en 6×6); donde no
cierra, al menos 71 veces más (4 lotes en 6×6, que el Split resuelve en 8,4
segundos).

La razón es estructural, como anticipaba el
Cuadro~\in[tab:diferencias_modelos]. El \ac{mip} del monolítico crece con la
grilla, porque incluye el ruteo de todos los lotes, mientras que el de la
Etapa~2 no depende de ella, porque recibe las rutas ya resueltas
(Cuadro~\in[tab:res_tamano_mip]). La Etapa~1, por su parte, consiste en
\ac{mip} chicos de ruteo, uno por lote, que en estas instancias se resuelven
en pocos segundos.

\startplacetable[reference=tab:res_tamano_mip, location=here, title={Tamaño del \ac{mip} que resuelve SCIP: el modelo monolítico completo y la Etapa~2 del Split.}]
	\bTABLE[split=yes, align=middle, frame=on, width=fit]
	\setupTABLE[row][1][background=color, backgroundcolor=darkgray, style=bold]
	\bTR \bTH Instancia \eTH \bTH Monolítico: variables \eTH \bTH Monolítico: restricciones \eTH \bTH Etapa~2: variables \eTH \bTH Etapa~2: restricciones \eTH \eTR
	\bTR \bTD 3 lotes, 3×3 \eTD \bTD 327 \eTD \bTD 573 \eTD \bTD 174 \eTD \bTD 339 \eTD \eTR
	\bTR \bTD 3 lotes, 5×5 \eTD \bTD 663 \eTD \bTD 1.029 \eTD \bTD 174 \eTD \bTD 339 \eTD \eTR
	\bTR \bTD 3 lotes, 6×6 \eTD \bTD 903 \eTD \bTD 1.347 \eTD \bTD 174 \eTD \bTD 339 \eTD \eTR
	\bTR \bTD 3 lotes, 7×7 \eTD \bTD 1.191 \eTD \bTD 1.725 \eTD \bTD 174 \eTD \bTD 339 \eTD \eTR
	\bTR \bTD 4 lotes, 3×3 \eTD \bTD 438 \eTD \bTD 772 \eTD \bTD 234 \eTD \bTD 460 \eTD \eTR
	\bTR \bTD 4 lotes, 5×5 \eTD \bTD 886 \eTD \bTD 1.380 \eTD \bTD 234 \eTD \bTD 460 \eTD \eTR
	\bTR \bTD 4 lotes, 6×6 \eTD \bTD 1.206 \eTD \bTD 1.804 \eTD \bTD 234 \eTD \bTD 460 \eTD \eTR
	\eTABLE
\stopplacetable

Más tiempo no resuelve el problema. Con un tope de 40 minutos, el monolítico
sigue sin probar el óptimo en esos puntos, ni en la instancia del escenario
de estrés 04 resuelta con el peso de finalización de la exploración
($W^{c} = 9$), como muestra el Cuadro~\in[tab:res_40min]. Encuentra
soluciones muy buenas, a menos del 1~\% del óptimo, pero no puede demostrar
que lo son: su cota dual casi no avanza, y en la instancia del escenario 04
informa un gap de 143,5~\% al cabo de 40 minutos, cuando el Split prueba el
óptimo en 181 segundos. Sin el óptimo del Split, la solución del monolítico
no podría distinguirse de una mala.

\startplacetable[reference=tab:res_40min, location=here, title={Resultado del monolítico con un tope de 40 minutos, frente al tiempo del Split en la misma instancia. La distancia al óptimo se mide contra el óptimo del Split.}]
	\bTABLE[split=yes, align=middle, frame=on, width=fit]
	\setupTABLE[row][1][background=color, backgroundcolor=darkgray, style=bold]
	\bTR \bTH Instancia \eTH \bTH Split (s) \eTH \bTH Monolítico \eTH \bTH ¿Prueba el óptimo? \eTH \bTH Distancia al óptimo \eTH \bTH Gap informado \eTH \eTR
	\bTR \bTD 3 lotes, 7×7 \eTD \bTD 8,2 \eTD \bTD 40,0 min \eTD \bTD no \eTD \bTD 0,05~\% \eTD \bTD 11,0~\% \eTD \eTR
	\bTR \bTD 4 lotes, 6×6 \eTD \bTD 9,5 \eTD \bTD 40,0 min \eTD \bTD no \eTD \bTD 0,00~\% \eTD \bTD 10,4~\% \eTD \eTR
	\bTR \bTD 6 lotes, 5×5 \eTD \bTD 231,7 \eTD \bTD 32,0 min \eTD \bTD sí \eTD \bTD 0,00~\% \eTD \bTD -- \eTD \eTR
	\bTR \bTD Estrés 04 ($W^{c} = 9$): 6 lotes, 7×7 \eTD \bTD 181,4 \eTD \bTD 40,0 min \eTD \bTD no \eTD \bTD 0,66~\% \eTD \bTD 143,5~\% \eTD \eTR
	\eTABLE
\stopplacetable

Estos estudios se ejecutaron con la máquina compartida con otras corridas,
por lo que sus tiempos absolutos están inflados. Lo que se compara es la
relación entre los dos modelos sobre la misma instancia, resueltos uno
después del otro.

\stopsubsection

\startsubsection[title={Pruebas de estrés}]

En las pruebas de estrés, el monolítico no termina ninguno de los 19
escenarios dentro del tope de 40 minutos; en el 01, el proceso terminó sin
devolver un estado. El Split termina 9, ninguno en más de 15,1 minutos. Los
escenarios que el Split tampoco termina son los del perfil realista con 8
lotes o más (12 a 14 y 16 a 19), la grilla de 16×16 (08) y dos de los
heredados (01 y 03). Son dos límites distintos: la Etapa~2 se encarece con
la cantidad de lotes y la Etapa~1, con el tamaño del almacén, como ya
mostraba el relevamiento (Sección~\in[sec:relevamiento]).

En síntesis, las dos formulaciones alcanzan el mismo óptimo, pero solo el
Split lo prueba a escala realista. Por eso, el Split es la referencia contra
la que se comparan las metaheurísticas, y el monolítico queda como juez de la
verificación de sus soluciones (Sección~\in[sec:validacion]).

\stopsubsection

\stopsection

%----------------------------------------------------------------------------------------
\startsection[title={Comparación entre modelo Split y metaheurísticas}, reference=sec:comparacion_split_metaheuristicas]

% TODO: completar

\stopsection
\stopchapter

%========================================================================================
% GUÍA DE TRABAJO - CAPÍTULO 6 (Resultados)
% Comentarios: no se compilan ni salen en el PDF. Borrar antes de la entrega.
% Armada el 21/10/2026 contrastando el capítulo con el repo de implementación
% (linear-picking-routing, origin/main 75d553c del 19/10), docs/thesis/models.tex, las dos
% comparaciones de docs/thesis/, los CSV de results/ (recalculados el 21/10) y
% Notas/opp_v5_development.tex. Foco: comparación entre modelos y validación de las
% metaheurísticas. Los números de abajo salen de las corridas consolidadas actuales.
% [ ] = tarea; (D) = decisión de ustedes.
%========================================================================================
%
% 0. ANTES DE ESCRIBIR NÚMEROS
% (D) Versión de las heurísticas. Las corridas consolidadas de HC y ACO (funcional
%     test_v5eq__* y estrés) son del 28 y 29/09, anteriores al commit 75d553c (19/10), que
%     fijó el orden de la ruleta de las construcciones aleatorias. Con el código actual la
%     misma semilla ya no da esas soluciones, y antes dependían de PYTHONHASHSEED (según el
%     commit, el 14/10: hasta 0,5 % de dispersión en el 03 y el 11 infactible en 2 de 7
%     procesos). Opciones: (a) volver a correr HC y ACO (funcional y estrés, con
%     WRITE_RESULTS_CSV=1 y WRITE_SOLUTIONS=1), verificarlos con run_compliance, regenerar el
%     registro (python -m analysis.comparability) y el documento
%     (experiments/opp_v5/reports/build_heuristics_comparison.py). Por los tiempos actuales,
%     unos 20 min HC (casi todo el 100×100) y 1,5 h ACO (35 min el 100×100), más unos
%     minutos de verificación; con la notebook enchufada y sin suspensión. (b) Dejar los
%     números del 28-29/09 y declarar la versión del código y esa dependencia.
% (D) Tiempos de estrés del Split con la Etapa 1 en 6 procesos (ver la guía del cap. 5).
% [ ] No usar analysis/output/report.md: su sección 4.1 mezcla tests no equivalentes (brechas
%     de hasta 1.180 %). Las fuentes son results/comparability/*.csv y los documentos
%     generados de docs/thesis/ (models.tex y las dos comparaciones), que salen de los CSV
%     consolidados.
%
% 1. ORGANIZACIÓN PROPUESTA (las dos secciones actuales pasan a 6.1 y 6.3)
% 6.1 Comparación entre modelos lineales
%   6.1.1 Mismo óptimo: 49 de 49 pares resueltos con el mismo objetivo (diferencia máxima
%         3,83 · 10^-6, casos 13 y 14) y 11 de 11 infactibles en los dos; 28 con la misma
%         solución y 21 óptimos alternativos (13 con otras rutas, 16 con otra asignación).
%         Explicar por qué: el objetivo no penaliza la longitud de la ruta, y los empates de
%         asignación los rompe cada resolución a su manera.
%   6.1.2 Instancias chicas (funcional): suma de total_time 140,3 contra 24,4 s (5,75×), pero
%         mediana 0,225 contra 0,309 s: el monolítico es más rápido en 38 de 49 pares
%         (compila un MIP en vez de dos) y hasta 28,9 veces más lento en las grillas (caso
%         29); model_time 17,35 veces en la suma. Cuadro y gráfico de dispersión.
%   6.1.3 Escala (estudios 08 a 11, la misma instancia en los dos modelos): frontera a 10 min
%         (3 lotes en 6×6, 4 en 5×5 y 6 en 3×3; el Split cierra todos); el cociente crece
%         con la grilla (11 a 15 veces con 2 lotes, hasta 27 con 3 y 4 lotes, y al menos 59 y
%         68 donde se cortó); encontrar no es el problema, probar sí: soluciones a menos de
%         1 % antes de 75 s, pero la cota dual se estanca entre el 88 % y el 97 % del óptimo;
%         40 min no lo arreglan (estrés 04 con W^c = 9: solución a 0,66 % y gap de 143,5 %).
%         Tamaño de los MIP (cuadro del cap. 3). Decidir si la figura de la frontera queda acá
%         o en 5.3.1.4.
%   6.1.4 Estrés: Split 9 de 19, monolítico 0 de 19. Cuadro por test con total, Etapa 1
%         estimada (total menos Etapa 2) y Etapa 2: domina la Etapa 1 salvo en el 04 (7×7) y
%         en los de agenda pesada (10 y 15); de 13×13 a 14×14 se multiplica por 2,0 y de
%         14×14 a 15×15 por 4,6; el 8×21 tarda más que el 14×14 en la Etapa 1 (pesa la
%         forma); el Split no termina con 8 lotes o más ni en 16×16: dos paredes (la Etapa 2
%         con los lotes, la Etapa 1 con la grilla).
%   6.1.5 Cierre: el Split es la referencia de las heurísticas; el monolítico, la de la
%         verificación.
% 6.2 Validez de las soluciones de las metaheurísticas (sección nueva; el método está en el
%     cap. 5)
%   - 150 de 150 "compliant": 56 funcionales y 19 de estrés por heurística, con el objetivo
%     recalculado igual al reportado en todas. Además dan compliant las soluciones de las
%     suites funcionales propias (51 por heurística) y las 9 de estrés del Split verificadas
%     contra el monolítico. Las infactibles (12 por heurística) no tienen solución que
%     verificar.
%   - Donde el modelo exacto esperaría (y en el 30a), el nivel 1 da un objetivo menor que el
%     reportado, hasta un 96 % menos (26b): es la espera voluntaria con las mismas
%     decisiones, no una violación.
%   - Notas del verificador: lotes vacíos quitados en 4 tests, M reemplazada en 1 (de la
%     suite propia) y subgrafo en el 100×100 (HC: 1.425 de 10.000 nodos y 2.900 de 39.600
%     arcos; ACO: 1.499 nodos y 3.160 arcos).
%   - Ninguna brecha negativa: ninguna heurística queda por debajo de un óptimo probado.
% 6.3 Split contra metaheurísticas
%   6.3.1 Estado: 67 de 67 funcionales comparables y 9 de 9 de estrés, con la lectura
%         correcta (no inventan soluciones donde no las hay; no prueban infactibilidad).
%         Desglosar las 12 infactibles por causa (cap. 5): las de validación de entrada se
%         detectan con la misma regla que en los exactos; las de ruta, por búsqueda con tope;
%         las de frío, por el prechequeo o por agotar la búsqueda.
%   6.3.2 Calidad: funcional, ACO 45 de 45 óptimas; HC 40 de 45 (media 0,36 %, máxima
%         5,88 %; las 5 que falla son de grilla: 20, 25c, 25e·1, 25f·1 y 25f·2). Estrés: ACO
%         2 de 9 óptimas, media 0,01 % y máxima 0,06 %; HC 1 de 9, media 0,06 % y máxima
%         0,21 % (el 11, con cadena de frío); ACO mejor que HC en 8 de 9. Gráfico de brecha
%         por escenario.
%   6.3.3 Instancias con espera voluntaria (10): cuadro Split / HC / ACO con los lotes que
%         esperan; se comparan solo en estado y la diferencia no es una brecha (en el 26b el
%         objetivo del Split es unas 30 veces menor; en las demás, entre 1,3 y 13 veces).
%   6.3.4 Tiempos: funcional, sobre las 33 instancias con tiempo del Split conocido (su CSV
%         guarda una fila por método): medianas de 0,300 / 0,061 / 0,224 s, cocientes de 5,1
%         (HC) y 1,4 (ACO); ACO es más lento que el Split en 9 de las 33. Estrés: cociente
%         mediano de 53,1 (HC) y 1,7 (ACO); HC en menos de 5 s salvo el 100×100 (17 min); ACO
%         de 13 s a 10 min (35 min en el 100×100). Gráficos: tiempo por escenario y tiempo
%         contra lotes × ubicaciones.
%   6.3.5 Sin referencia (los 10 escenarios que el Split no termina): las dos heurísticas
%         dan solución válida; ACO es mejor en 9 de 10 (por menos de 1 %); en el 01 ACO hizo
%         2 iteraciones y la notebook se suspendió: vale como solución factible, no como
%         medida de calidad.
%   6.3.6 A igual tiempo: con el tiempo de HC, ningún modelo exacto obtiene solución en los
%         19; con el de ACO, el Split da solución en 15 de 19 y ACO empata o mejora en 14
%         (en el 04 el Split es mejor por 0,001 %): en 8 a 0,1 % o menos, en el 02 y el 03 a
%         menos de 1 %, y entre 30 y 62 % a favor de ACO con 12 lotes o más (14, 16-19); sin
%         solución en el 01, 08 y 09 (la Etapa 1 no rutea todos los lotes) y en el 11 (rutas
%         cortadas que dejan lotes sin picker elegible por frío); el monolítico solo en el 04
%         y el 12 (entre 5 y 10 % peor que ACO). Advertencias en el cap. 5.
%   6.3.7 Nivel de servicio, memoria y convergencia (o a un anexo): SLA 100 % en el perfil
%         realista salvo el 11 (Split 97,4; HC 96,8; ACO 97,4) y 0 % en 01-03; memoria
%         mediana de 27 MB (HC) y 28 MB (ACO), 867 y 871 MB en el 100×100, y picos de ACO de
%         hasta 1.658 MB en 06, 09 y 12 sin explicar; HC corre entre 60 y 200 iteraciones (0
%         en el 01) y su solución inicial ya es la mejor en 3 escenarios; ACO corre las 60
%         iteraciones y su mejor solución aparece entre la 1 y la 54.
%   6.3.8 Variabilidad entre semillas: hoy no hay datos con el código final (una corrida
%         con semilla 1; los CSV de results/*/stability son del 11/09, con otras versiones
%         de los modelos y otras instancias). Ver 2.
% 6.4 Síntesis: quién domina en qué (gráfico de brecha contra tiempo) y la recomendación, o
%     dejarla para el cap. 7.
%
% 2. QUÉ FALTA
% [ ] Variabilidad entre semillas (quedó para después el 29/09): con el código final, p. ej.
%     10 semillas en los 19 escenarios de estrés (HC es barato; en ACO dejar afuera el 01),
%     reportando brecha mínima, mediana y máxima, y en cuántas semillas se encuentra
%     solución en el 11. Es necesario si "robustez" queda en los objetivos. El estudio de
%     cinco sorteos de HC del cap. 4 puede venir acá.
% [ ] Responder explícitamente los objetivos específicos del cap. 1 (convergencia, tiempo,
%     calidad, robustez y limitaciones estructurales), al final de cada sección o en 6.4.
% [ ] Sensibilidad a los parámetros de las heurísticas, si se muda del cap. 4: HC 200/60
%     contra 500/120; ACO, ablación y feromonas (results/opp_aco_4/{ablation,pheromone}).
%
% 3. QUÉ CORREGIR O MEJORAR (presentación)
% [ ] Colores: la guía de estilo asigna darkblue a los modelos exactos, blueline a HC y
%     greenline a ACO; los gráficos de models.tex y de las comparaciones usan greenline para
%     HC y redline para ACO. Unificar al reutilizarlos (el cap. 5 ya usa darkblue y redline
%     para las etapas del Split).
% [ ] Nombres: "monolítico", "Split", "HC" y "ACO" (no opp_v5_*, HC-4 ni ACO-4).
% [ ] Cada tiempo con qué mide (total o solver) y en qué condiciones (una corrida, una
%     máquina, Etapa 1 en paralelo, suspensiones); brecha y cociente con su fórmula.
% [ ] Conteos funcionales coherentes en todo el texto: 70 instancias = 68 que corren las
%     heurísticas + el 36c (no se captura) + el 37a (excluido); 64 métodos en el Split y 63
%     en el monolítico.
%
% 4. FUENTES
% - linear-picking-routing/docs/thesis/: models.tex (sec. "Comparación de los modelos
%   exactos con las heurísticas", con figuras MetaPost listas para reutilizar),
%   opp_v5_1_mono_vs_opp_v5_2_split_comparison.tex,
%   opp_v5_2_split_vs_heuristics_comparison.tex y opp_v5_1_mono_exploration.tex (estudios
%   08 a 11).
% - results/comparability/{comparability_v5_vs_heuristics, results_v5_vs_heuristics,
%   compare_v5_split_vs_mono}.csv; results/compliance/opp_{hc,aco}_4/; results/equal_time/
%   (27/09); los CSV de nombre fijo de results/opp_v5_*/ y results/opp_{hc,aco}_4/. Si
%   cambian, regenerar con experiments/opp_v5/reports/build_*.py.
% - Notas/opp_v5_development.tex: secs. "Resultados de las suites" y "El Split como
%   referencia para las heurísticas".
%========================================================================================
